%% ============================================================
%%  PRISM ISM 数据库 — 用户手册（中文版）
%%  编译方式：xelatex prism_manual_zh.tex（运行两次以生成目录）
%% ============================================================
\documentclass[12pt,a4paper]{article}

% ---- 中文支持 ----
\usepackage[UTF8, fontset=windows]{ctex}

% ---- 页面布局 ----
\usepackage[a4paper, top=2.5cm, bottom=2.5cm, left=2.8cm, right=2.8cm]{geometry}
\usepackage{setspace}
\onehalfspacing

% ---- 颜色 ----
\usepackage[dvipsnames,svgnames,table]{xcolor}
\definecolor{prismblue}{HTML}{4f6f9e}
\definecolor{prismclay}{HTML}{bf6a4a}
\definecolor{prismgreen}{HTML}{3f8a5e}
\definecolor{prismamber}{HTML}{c79a3e}
\definecolor{prismpurple}{HTML}{7d5ba6}
\definecolor{lightgray}{HTML}{f4f5f7}
\definecolor{midgray}{HTML}{e2e5ea}
\definecolor{darkink}{HTML}{2b3344}

% ---- 超链接 ----
\usepackage[colorlinks=true, linkcolor=prismblue, urlcolor=prismblue, citecolor=prismblue,
            pdftitle={PRISM ISM 数据库 — 用户手册},
            pdfauthor={PRISM Platform},
            pdfsubject={外资审查机制数据平台}]{hyperref}

% ---- 图形与文本框 ----
\usepackage{graphicx}
\usepackage{tcolorbox}
\tcbuselibrary{skins, breakable, theorems}
\usepackage{booktabs}
\usepackage{tabularx}
\usepackage{array}
\usepackage{longtable}
\usepackage{multirow}
\usepackage{enumitem}
\usepackage{fontawesome5}
\usepackage{mdframed}
\usepackage{fancyhdr}
\usepackage{titlesec}
\usepackage{titletoc}
\usepackage{parskip}

% ---- 章节格式 ----
\titleformat{\section}{\large\bfseries\color{darkink}}{}{0em}{\thesection\quad}[\vspace{-0.3em}\color{midgray}\rule{\linewidth}{0.8pt}]
\titleformat{\subsection}{\normalsize\bfseries\color{prismblue}}{}{0em}{\thesubsection\quad}
\titleformat{\subsubsection}{\normalsize\bfseries\color{darkink}}{}{0em}{\thesubsubsection\quad}

% ---- 页眉页脚 ----
\pagestyle{fancy}
\fancyhf{}
\fancyhead[L]{\small\color{gray}\textbf{PRISM} ISM 数据库 — 用户手册}
\fancyhead[R]{\small\color{gray}\leftmark}
\fancyfoot[C]{\small\color{gray}\thepage}
\renewcommand{\headrulewidth}{0.4pt}
\renewcommand{\footrulewidth}{0pt}

% ---- 自定义文本框 ----
\newtcolorbox{tipbox}[1][]{
  enhanced, breakable,
  colback=prismblue!6, colframe=prismblue!60,
  fonttitle=\bfseries\small, title={#1},
  left=6pt, right=6pt, top=5pt, bottom=5pt,
  arc=4pt
}
\newtcolorbox{warnbox}[1][]{
  enhanced, breakable,
  colback=prismclay!7, colframe=prismclay!60,
  fonttitle=\bfseries\small, title={#1},
  left=6pt, right=6pt, top=5pt, bottom=5pt,
  arc=4pt
}
\newtcolorbox{notebox}[1][]{
  enhanced, breakable,
  colback=lightgray, colframe=midgray,
  fonttitle=\bfseries\small, title={#1},
  left=6pt, right=6pt, top=5pt, bottom=5pt,
  arc=4pt
}
\newtcolorbox{keybox}{
  enhanced,
  colback=prismgreen!6, colframe=prismgreen!50,
  left=6pt, right=6pt, top=4pt, bottom=4pt,
  arc=4pt
}

% ---- 其他 ----
\usepackage{calc}
\setlist[itemize]{itemsep=2pt, topsep=4pt}
\setlist[enumerate]{itemsep=2pt, topsep=4pt}

% ---- 标签辅助命令 ----
\newcommand{\badge}[2]{%
  \colorbox{#1}{\color{white}\footnotesize\textbf{\strut\,#2\,}}%
}
\newcommand{\keycap}[1]{%
  \texttt{\footnotesize[#1]}%
}
\newcommand{\tabname}[1]{\textbf{\color{prismblue}#1}}
\newcommand{\uielem}[1]{\textit{#1}}
\newcommand{\fieldname}[1]{\texttt{\small#1}}

%% ============================================================
\begin{document}

% ---- 封面 ----
\begin{titlepage}
  \centering
  \vspace*{3cm}

  {\Huge\bfseries\color{prismclay} PRISM}\\[0.4em]
  {\LARGE\bfseries\color{darkink} ISM 数据库}\\[0.6em]
  {\large\color{gray} 外国直接投资审查机制数据平台}

  \vspace{1.5cm}
  \color{midgray}\rule{0.5\linewidth}{1pt}\color{black}
  \vspace{1.5cm}

  {\Large\bfseries 用户手册}\\[0.8em]
  {\normalsize 版本 1.0 \quad|\quad 数据范围：2007--2023}\\[0.4em]
  {\normalsize \url{https://jonebroom.github.io/prism-fdi/}}

  \vspace{2cm}
  \begin{keybox}
    \centering
    \textbf{38 个国家}\quad $\cdot$\quad \textbf{17 年}\quad $\cdot$\quad \textbf{6 个分析视图}\\
    \small OECD · 欧盟/欧洲经济区 · 五眼联盟 · 650 条时间序列记录 · 143 部法规 · 141 份政策文件
  \end{keybox}

  \vfill
  {\small\color{gray} 本手册涵盖平台全部功能、数据字段及分析模块。\\
  面向法律从业者、投资专业人士、合规团队及学术研究人员。}
\end{titlepage}

% ---- 目录 ----
\tableofcontents
\newpage

%% ============================================================
\section{平台概述}
\label{sec:overview}

\textbf{PRISM}（政策审查与投资筛查机制数据库，Policy Review and Investment Screening Mechanisms）是一个交互式数据平台，提供 2007 至 2023 年间 38 个国家外国直接投资（FDI）审查机制的全面时间序列数据。

\subsection{什么是 FDI 审查？}

外国直接投资审查是指政府对入境外资进行审核，并在国家安全、公共秩序或战略经济利益方面对交易加以限制、设定条件甚至予以否决的制度程序。自 2010 年代中期以来，随着地缘政治紧张局势加剧、新冠疫情冲击以及欧盟 2019 年《FDI 审查条例》出台，建立正式审查机制的国家数量急剧增加。

\subsection{平台适用对象}

\begin{itemize}
  \item \textbf{法律从业者与合规团队}——快速判断拟议交易是否触发目标司法管辖区的强制申报义务。
  \item \textbf{投资专业人士}——跨潜在目标市场比较审查制度，评估交易时间线风险。
  \item \textbf{政策研究者与学者}——追踪 FDI 审查规范的扩散与收紧趋势。
  \item \textbf{政府官员}——将本国机制与同类国家进行基准比较。
\end{itemize}

\subsection{数据范围}

\begin{table}[h]
\centering
\begin{tabularx}{\linewidth}{lX}
\toprule
\textbf{维度} & \textbf{覆盖范围} \\
\midrule
国家数量 & 38 个（OECD 成员国、欧盟/欧洲经济区国家、五眼联盟） \\
时间范围 & 2007--2023 年（年度观测） \\
记录数量 & 650 条国家-年度时间序列条目 \\
法规数量 & 143 部具体监管文件 \\
立法事件 & 230 次立法变更事件 \\
政策文件 & 141 份全文来源文件（PDF + HTML） \\
追踪行业 & 36 个细分行业 + 8 个聚合类别 \\
程序性指标 & 13 个二元权力指标 + 4 个定量字段 \\
\bottomrule
\end{tabularx}
\caption{PRISM 数据覆盖摘要}
\end{table}

\begin{tipbox}[数据时效性]
所有数据反映截至 2023 年末的已发布信息。2024 年数据因覆盖不完整而被刻意排除。平台默认显示最近完整年度（2023 年）的数据。
\end{tipbox}

\subsection{国家分组}

平台中国家被划分为三个相互重叠的分组，用于筛选过滤与颜色区分：

\begin{itemize}
  \item \badge{prismblue}{OECD}——经济合作与发展组织成员国
  \item \badge{prismgreen}{欧盟/欧洲经济区}——欧洲联盟及欧洲经济区成员国
  \item \badge{prismamber}{五眼联盟}——澳大利亚、加拿大、新西兰、英国、美国
\end{itemize}

一个国家可同时属于多个分组（例如澳大利亚既属于 OECD，也属于五眼联盟）。

%% ============================================================
\section{界面布局}
\label{sec:layout}

PRISM 界面分为四个始终可见的固定区域，以及按需弹出的侧边抽屉。

\begin{figure}[h]
\centering
\begin{tcolorbox}[enhanced, colback=lightgray, colframe=midgray, arc=6pt, width=\linewidth]
\small
\begin{verbatim}
┌─────────────────────────────────────────────────────────────┐
│  顶栏：Logo │ 搜索 │ ⚡ 交易核查 │ AI 助手 │ ℹ 信息        │
├──────────┬──────────────────────────────────────────────────┤
│          │  标签栏：01 概览 │ 02 国家 │ … │ 06 行业        │
│  侧边栏  ├──────────────────────────────────────────────────┤
│  年份    │                                                   │
│  分组    │              主内容区域                          │
│  国家    │            （当前标签面板）                      │
│  列表    │                                                   │
└──────────┴──────────────────────────────────────────────────┘
\end{verbatim}
\end{tcolorbox}
\caption{PRISM 界面布局示意图}
\end{figure}

\subsection{顶部导航栏}

顶部导航栏横贯全宽，包含以下元素：

\begin{itemize}
  \item \uielem{PRISM 标志}——平台标识，副标题显示国家数量与年份范围。
  \item \uielem{全局搜索框}——跨国家、法规及政策文件的全文搜索。在任意位置按 \keycap{/} 键可聚焦搜索框。
  \item \uielem{⚡ 交易核查}——打开交易筛查抽屉（详见第 \ref{sec:screener} 节）。
  \item \uielem{AI 助手}——打开 AI 数据助手抽屉（详见第 \ref{sec:ai} 节）。
  \item \uielem{ℹ}——打开平台指南抽屉，提供快速参考信息。
\end{itemize}

\subsection{控制侧边栏}
\label{sec:rail}

左侧边栏包含影响整个平台的全局控件：

\subsubsection{年份选择器}

\begin{itemize}
  \item 大号年份数字显示当前选定年份。
  \item 左右拖动\uielem{滑块}在 2007 年至 2023 年之间切换。
  \item 点击\uielem{▶ 播放时间轴}可自动逐年动态演示（每秒一年）；点击\uielem{❚❚ 暂停}停止。切换至其他标签页也会停止动画。
\end{itemize}

\subsubsection{国家分组筛选器}

三个切换按钮用于控制哪些国家出现在所有图表和列表中：

\begin{itemize}
  \item \badge{prismblue}{OECD}（19 个国家）——点击以启用/禁用。
  \item \badge{prismgreen}{欧盟/欧洲经济区}（27 个国家）——点击以启用/禁用。
  \item \badge{prismamber}{五眼联盟}（5 个国家）——点击以启用/禁用。
\end{itemize}

\begin{warnbox}[注意]
至少须保持一个分组处于激活状态。若三个分组均被停用，最后一个切换的分组将自动重新启用。由于各分组存在国家重叠，当前激活国家总数可能少于各分组规模之和。
\end{warnbox}

\subsubsection{国家列表}

显示通过当前分组筛选的所有国家，每行显示：
\begin{itemize}
  \item 彩色圆点，表示所属分组。
  \item 国家名称。
  \item 审查机制建立年份（若无则显示"—"）。
\end{itemize}

使用列表上方的\uielem{筛选国家……}搜索框可快速定位。点击任意国家可将其选中，并在列表中高亮显示，同时自动切换至\tabname{国家}标签页以展示该国详情。

\subsection{标签栏}

平台提供六个分析标签页。每个标签页独立运行，但均响应侧边栏中全局年份与国家选择。

\begin{table}[h]
\centering
\begin{tabularx}{\linewidth}{clX}
\toprule
\textbf{编号} & \textbf{标签名称} & \textbf{功能} \\
\midrule
01 & 概览 & 世界地图 + 全球趋势折线图 + 年度快照 \\
02 & 国家 & 所选国家的深度数据档案 \\
03 & 比较 & 多国平行坐标图 \\
04 & 演变 & 覆盖类型趋势、严格程度分布、监管网络 \\
05 & 变化 & 按年度的全球立法活动 \\
06 & 行业 & 行业覆盖热图与排名 \\
\bottomrule
\end{tabularx}
\caption{六个分析标签页}
\end{table}

\subsection{移动端使用}

在屏幕宽度小于 768\,像素的设备（智能手机及小型平板）上：
\begin{itemize}
  \item 侧边栏默认隐藏，点击左上角\uielem{☰}图标可打开。
  \item 点击侧边栏外部深色遮罩或\uielem{✕}按钮可关闭侧边栏。
  \item 选择国家后侧边栏自动关闭。
  \item 标签栏变为可横向滚动。
  \item 所有图表面板纵向堆叠显示。
\end{itemize}

\subsection{键盘快捷键}

\begin{table}[h]
\centering
\begin{tabular}{ll}
\toprule
\textbf{按键} & \textbf{操作} \\
\midrule
\keycap{/} & 聚焦全局搜索框 \\
\keycap{Esc} & 关闭任意打开的抽屉 \\
\keycap{Enter}（在搜索框中） & 跳转至第一条结果 \\
\keycap{↑}/\keycap{↓}（在搜索框中） & 在结果列表中上下移动 \\
\bottomrule
\end{tabular}
\caption{键盘快捷键}
\end{table}

%% ============================================================
\section{标签页 01 — 概览}
\label{sec:overview-tab}

概览标签页是平台的起始视图，提供所选年份全球 FDI 审查活动的综合快照。

\subsection{世界地图}

\begin{itemize}
  \item 每个有数据的国家根据所选\uielem{维度}（在面板标题的下拉菜单中选择）进行深浅着色。
  \item 颜色越深，数值越高；不在当前分组筛选中的国家以浅灰色显示。
  \item \textbf{可用地图维度：}
    \begin{itemize}
      \item \fieldname{严格程度评分（0--13）}——13 项程序性权力的总和。
      \item \fieldname{覆盖行业数}——纳入审查机制的行业数量。
      \item \fieldname{审查触发门槛（\%）}——触发审查的最低股权比例。
      \item \fieldname{审查时限（天）}——法定最长审查期。
    \end{itemize}
  \item \textbf{点击}地图上任意国家可弹出深色浮动卡片，显示关键指标。点击\uielem{查看详情 →}可跳转至该国完整档案。
  \item 地图支持\textbf{缩放与平移}：滚动鼠标滚轮缩放，拖动平移；双击可重置视图。
\end{itemize}

\subsection{全球趋势图}

下方左侧折线图追踪每年有特定功能处于激活状态的国家数量，叠加显示每年全球新建审查机制数量的柱状图。

\begin{itemize}
  \item 使用\uielem{下拉菜单}更改追踪功能（如国家安全测试、罚款、缓解措施等）。
  \item 两条竖向虚线分别标记 \textbf{2019 年}（欧盟 FDI 法规推动，蓝色）和 \textbf{2020 年}（新冠疫情应急立法浪潮，橙色）。
  \item \textbf{点击}图表上任意年份可将全局年份滑块跳转至该年。
\end{itemize}

\subsection{年度快照面板}

右侧面板显示当前激活国家集合在所选年份的汇总数据：

\begin{itemize}
  \item \textbf{已建立审查机制的国家数}——筛选国家中拥有有效审查制度的数量。
  \item \textbf{平均严格程度}——所有有审查机制国家的平均严格程度评分。
  \item \textbf{平均覆盖行业数}——平均行业覆盖数量。
  \item \textbf{本年新增机制数}——全球新建审查机制数量。
  \item \textbf{覆盖类型分布}——横向条形图，显示各覆盖模式（行业性、跨行业、混合、资产型等）对应的国家数量。
\end{itemize}

%% ============================================================
\section{标签页 02 — 国家}
\label{sec:country}

国家标签页提供单个国家 FDI 审查机制的最详尽视图。可从侧边栏列表或点击世界地图选择国家。

\subsection{从业者摘要卡片}

顶部卡片专为快速实务参考而设计：

\subsubsection{状态标签}

卡片右上角的彩色标签汇总整体风险等级：

\begin{table}[h]
\centering
\begin{tabular}{lll}
\toprule
\textbf{标签} & \textbf{颜色} & \textbf{含义} \\
\midrule
\badge{gray}{无审查机制} & 灰色 & 所选年份无有效审查机制 \\
\badge{prismgreen}{审查机制存在} & 绿色 & 机制存在但严格程度较低（0--3 分） \\
\badge{prismamber}{主动审查} & 琥珀色 & 中等严格程度（4--7 分） \\
\badge{prismclay}{严格审查} & 橙红色 & 高严格程度（8--13 分） \\
\bottomrule
\end{tabular}
\caption{国家状态标签定义}
\end{table}

\subsubsection{申报要求模块}

回答"我是否需要申报？"这一实务问题：

\begin{itemize}
  \item \textbf{交割前须取得预先批准}——强制预批；未获政府批准，交易不得完成交割。
  \item \textbf{须强制通知}——必须申报，但交割可在决定作出前进行。
  \item \textbf{无强制申报}——仅需自愿通知，或无审查机制。
  \item \textbf{触发门槛}——触发审查义务的最低股权比例（\%）。
  \item \textbf{最长审查时限}——法定最长审查期（天）。
  \item \textbf{通知要求}——来源数据中具体的通知制度说明。
\end{itemize}

\subsubsection{审查机构模块}

\begin{itemize}
  \item \textbf{主管机构}——负责管理审查程序的政府机关。
  \item \textbf{审查标准}——可据以阻止投资的法律依据（国家安全、净经济效益、竞争、两用品/出口管制）。
  \item \textbf{严格程度评分}——13 项程序性权力中处于激活状态的数量。
  \item \textbf{覆盖行业数}——纳入审查范围的行业数量。
\end{itemize}

\subsubsection{范围与权力模块}

\begin{itemize}
  \item 是否涵盖绿地投资（新建项目，而非仅收购）。
  \item 是否涵盖房地产交易。
  \item 审查机制是否区分欧盟与非欧盟投资者。
  \item 功能标签，列出已激活权力：正式主动调查权、跨部门联合审查、分级管辖权、缓解措施权、罚款/处罚权、黄金股、欧盟与非欧盟投资者差异化处理。
\end{itemize}

\subsubsection{快速交易核查按钮}

卡片底部的\uielem{⚡ 针对该国核查具体交易}按钮将自动填入当前选定国家并立即打开交易筛查抽屉。

\subsection{多维雷达图}

在 8 个程序性维度上同时呈现所选国家的数据，并叠加两条参考线：

\begin{itemize}
  \item \textbf{橙红色实线}——所选国家当前年份数据。
  \item \textbf{灰色虚线}——当前年份 OECD 平均值。
  \item \textbf{蓝色点线}——所选国家历史对比年份数据（通过\uielem{对比 [年份]}下拉菜单选择）。
\end{itemize}

\textbf{雷达图维度：}国家安全测试、正式主动调查权、罚款、缓解措施、跨部门审查、净效益测试、强化政府控制权、股权审查。

\begin{tipbox}[雷达图使用技巧]
将对比年份切换至特定年份（如 2015 年），可直观看出该国审查机制的收紧幅度。使用全局年份滑块可审查特定历史时期的数据。
\end{tipbox}

\subsection{关键指标面板}

所选国家与年份的全部可用数据字段结构化列表，包括：主管机构、覆盖类型、通知制度、预批制度、审查门槛、审查时限、欧盟与非欧盟差异、新冠疫情临时立法标志。

\subsection{程序要求与覆盖范围}

16 个二元指标按四类分组显示：

\begin{table}[h]
\centering
\begin{tabularx}{\linewidth}{lX}
\toprule
\textbf{类别} & \textbf{指标} \\
\midrule
调查与审查权力 & 正式主动调查权、增量股权审查、跨部门联合审查、分级管辖权、强化政府控制权 \\
测试与标准 & 国家安全测试、净效益测试、竞争测试 \\
制裁与机制 & 罚款、缓解措施、申报费、当地代表、联合选址 \\
覆盖范围 & 绿地投资、房地产、出口管制/两用品 \\
\bottomrule
\end{tabularx}
\caption{按类别分组的程序性指标}
\end{table}

每项指标以 \textcolor{prismgreen}{\textbf{✓}}（存在）或 \textcolor{gray}{\textbf{✕}}（缺失）显示。

\subsection{法规时间轴}

水平气泡时间轴，显示该国所有已记录的监管文件：

\begin{itemize}
  \item 每个气泡的\textbf{大小}与该法规覆盖的行业数成正比。
  \item 每个气泡的\textbf{颜色}反映覆盖类型。
  \item \textbf{点击}任意气泡可打开详情抽屉，包含：
    \begin{itemize}
      \item 法规全称、签署年份/生效年份。
      \item 主管机构、通知制度、触发门槛、审查时限。
      \item 覆盖行业数。
      \item 替代信息（若该法规取代了早期法规）。
      \item 政策文本摘录（如有来源文件）。
      \item \uielem{用 AI 总结 →}按钮，将法规发送至 AI 助手。
      \item \uielem{↓ 下载 PDF} 或 \uielem{↗ 查看原文}链接。
    \end{itemize}
\end{itemize}

%% ============================================================
\section{标签页 03 — 比较}
\label{sec:compare}

比较标签页使用\textbf{平行坐标图}，在多个维度上同时呈现所有激活国家的数据。

\subsection{图表解读}

\begin{itemize}
  \item 每条\textbf{纵轴}代表一个指标。
  \item 每条\textbf{彩色线条}代表所选年份的一个国家。
  \item 线条在每条轴上的位置表示该国对应指标的数值。
  \item 无审查机制的国家不显示在图中。
\end{itemize}

\textbf{显示的坐标轴：}
机制数量、严格程度评分、覆盖行业数、触发门槛（\%）、审查时限（天）、是否有罚款、是否有缓解措施、是否有跨部门审查、是否有国家安全测试、是否有净效益测试。

\subsection{交互操作}

\begin{itemize}
  \item \textbf{在任意坐标轴上拖动}可创建筛选范围；只有通过选定区间的线条保持高亮，便于筛选（例如触发门槛低于 25\% 且有国家安全测试的所有国家）。
  \item \textbf{悬停}线条可高亮显示并查看包含所有数值的提示框。
  \item \textbf{点击}线条可跳转至该国在国家标签页中的完整档案。
  \item 使用\uielem{按分组着色 / 按严格程度着色}切换按钮，可在分组颜色编码与严格程度渐变着色之间切换。
\end{itemize}

\begin{tipbox}[典型使用场景]
选择目标国家并在严格程度轴上拖动至其数值，剩余高亮国家拥有相似的严格程度——便于同类比较或识别具有相近合规负担的司法管辖区。
\end{tipbox}

%% ============================================================
\section{标签页 04 — 演变}
\label{sec:evolution}

演变标签页包含五张图表，共同展现 FDI 审查制度随时间的结构性变化。

\subsection{覆盖类型分布（饼图）}

当前选定年份的饼图，显示各覆盖模式对应的国家数量：

\begin{itemize}
  \item \badge{prismclay}{行业性}——仅涵盖特定行业。
  \item \badge{prismblue}{跨行业}——所有行业均可能纳入审查范围。
  \item \badge{prismpurple}{混合型}——宽泛规则与行业专项规则相结合。
  \item \badge{prismamber}{资产型}——审查由资产类型触发，而非行业。
  \item \badge{gray}{草案}——立法已提出但尚未颁布。
  \item \badge{gray}{已通过（未实施）}——法律已通过，但实施细则或生效日期尚未到位。
\end{itemize}

\textbf{点击}任意扇区可打开抽屉，列出所选年份属于该类型的所有国家。

\subsection{通知 × 预批矩阵（气泡图）}

在两个坐标轴上绘制国家分布的 3×3 矩阵：

\begin{itemize}
  \item \textbf{X 轴（通知要求）：}非强制 · 强制（部分交易）· 强制（所有交易）
  \item \textbf{Y 轴（预批要求）：}非强制 · 强制（部分）· 强制（所有）
  \item \textbf{气泡大小：}该组合对应的国家数量。
\end{itemize}

\textbf{点击}任意气泡可查看具有该程序要求组合的国家列表。

\subsection{覆盖类型演变（面积图）}

堆积面积图，显示 2007 至 2023 年每年各覆盖模式的国家数量。通过此图可观察到 2019 年后从行业性机制向跨行业机制的结构性转变加速这一趋势。

\subsection{严格程度随时间分布（带状图）}

以折线绘制每年\textbf{平均严格程度评分}，并以带状区域显示\textbf{最小值和最大值}。带状区域的扩大反映了最严格与最宽松司法管辖区之间差距的不断扩大。

\subsection{门槛与严格程度 / 时限与严格程度（散点图）}

所选年份的散点图：
\begin{itemize}
  \item \textbf{X 轴}——审查触发门槛（\%）或审查时限（天），可通过\uielem{门槛/时限}切换按钮选择。
  \item \textbf{Y 轴}——严格程度评分（0--13）。
  \item 每个点代表一个国家。\textbf{点击}可跳转至该国档案。
\end{itemize}

位于左上方的国家（严格程度高、门槛低 / 时限短）对投资者施加了最严苛的实际合规负担。

\subsection{法规替代关系网络图}

有向网络图，其中：
\begin{itemize}
  \item \textbf{节点}为个别法规，按国家分组着色。
  \item \textbf{有向边}从旧法指向取代它的新法。
  \item 由边连接的节点集群代表某国的立法演变脉络。
  \item \textbf{点击}任意节点可导航至该国档案。
\end{itemize}

%% ============================================================
\section{标签页 05 — 变化}
\label{sec:changes}

变化标签页聚焦于立法活动——发生了什么，以及何时发生。

\subsection{全球立法活动柱状图}

堆积柱状图，显示每年全球发生的立法事件数量，并按类型分解：

\begin{itemize}
  \item \badge{prismblue}{新立法}——建立审查机制的全新法规。
  \item \badge{gray}{修正案}——对现有法律的修改。
  \item \badge{prismamber}{行政命令}——总统/行政令或紧急法令。
  \item \badge{gray}{实施条例}——此前颁布法律的配套实施细则。
\end{itemize}

\textbf{点击}任意柱体可选定该年份；下方事件列表将更新为该年所有事件，\textbf{全局年份滑块也会同步}至该年——切换至国家或概览标签页时将显示该年数据。

\subsection{立法事件列表}

每张事件卡片显示：
\begin{itemize}
  \item 年份与国家。
  \item 法律名称或引用编号（形如 \texttt{34/2020} 的裸引用编号标注"引用："，表示来源数据中无完整名称）。
  \item 变更内容简述（如有）。
  \item 事件类型标签：新立法、修正案、行政命令或实施条例。
  \item 若该立法取代了早期文件，则显示"替代"标签。
\end{itemize}

\textbf{点击}任意事件卡片可打开详情抽屉，包含：
\begin{itemize}
  \item 匹配机制数据（主管机构、触发门槛、严格程度、覆盖行业数），如能找到对应法规记录。
  \item 政策文本摘录及法规全文（如有）。
  \item \uielem{查看 [国家] 详情 →}按钮，跳转至国家标签页。
  \item \uielem{用 AI 总结 →}，向 AI 助手发起法规摘要请求。
\end{itemize}

%% ============================================================
\section{标签页 06 — 行业}
\label{sec:sectors}

行业标签页分析 FDI 审查跨国家、跨时间在不同行业中的覆盖情况。

\subsection{行业覆盖热图}

矩阵图，包含：
\begin{itemize}
  \item \textbf{行}——各具体行业（或 8 个聚合类别，可切换）。
  \item \textbf{列}——各国家，按行业覆盖总数排序（覆盖最多者居前）。
  \item \textbf{单元格颜色}——深色表示该行业已纳入审查；浅色表示未覆盖。
\end{itemize}

\textbf{两种显示模式：}
\begin{itemize}
  \item \uielem{细分行业}——显示所有 36 个具体行业名称（如"人工智能与机器学习"、"民用核能"、"量子技术"）。
  \item \uielem{8 大聚合类别}——以百分比（0--100\%）显示八个大类的覆盖率，采用渐变着色。
\end{itemize}

\subsection{行业覆盖排名（条形图）}

水平条形图，按将某行业纳入审查机制的国家数量对所有行业排名，覆盖国家越多的行业位置越靠前。

\begin{itemize}
  \item 图表反映当前所选年份的数据。
  \item 点击\uielem{▶ 播放}可逐年动态演示，观察随着新行业被纳入各国机制，排名的演变过程；点击\uielem{❚❚ 暂停}停止。
  \item 切换至其他标签页会自动停止动画。
\end{itemize}

\textbf{8 大行业聚合类别如下：}

\begin{table}[h]
\centering
\begin{tabularx}{\linewidth}{lX}
\toprule
\textbf{类别} & \textbf{典型子行业} \\
\midrule
国防 & 国防生产、受控两用品、出口管制 \\
实体关键基础设施 & 能源、水务、交通、医疗、电信 \\
新一代关键基础设施 & 航天、先进计算、量子技术、人工智能 \\
关键技术与两用品 & 网络安全、机器人技术、生物技术、先进材料 \\
非贸易性服务 & 博彩、旅游、教育、民用核能 \\
金融 & 银行、保险、支付 \\
媒体 & 广播、新闻、出版 \\
敏感个人数据 & 个人数据访问、脑机接口 \\
\bottomrule
\end{tabularx}
\caption{8 大行业聚合类别}
\end{table}

%% ============================================================
\section{交易筛查器}
\label{sec:screener}

交易筛查器回答一个核心实务问题：\textbf{"此次交易是否可能需要申报？"}

\begin{tipbox}[打开方式]
点击顶部导航栏中的\uielem{⚡ 交易核查}，或点击任意国家档案卡片上的\uielem{⚡ 核查具体交易}。
\end{tipbox}

\subsection{输入字段}

\begin{enumerate}
  \item \textbf{目标国家}——被投资企业所在国家。
  \item \textbf{行业/领域}——从 8 个大类或 36 个细分行业中选择；不确定可留空。
  \item \textbf{拟收购股权比例（\%）}——投资者拟收购的股权比例；不明确可留空。
  \item \textbf{您的国家}——投资者母国，用于在审查机制区分欧盟与非欧盟投资者时判断适用规则。
\end{enumerate}

\subsection{核查结论}

点击\uielem{核查此交易 →}后，筛查器返回五个结论等级之一：

\begin{table}[h]
\centering
\begin{tabularx}{\linewidth}{llX}
\toprule
\textbf{结论} & \textbf{颜色} & \textbf{含义} \\
\midrule
可能须接受审查 & \textcolor{prismclay}{橙红色} & 行业已涵盖、股权比例超过门槛、须强制预批 \\
可能适用审查 & \textcolor{prismamber}{琥珀色} & 部分因素提示须申报，但确定性有限 \\
低于门槛 & \textcolor{prismgreen}{绿色} & 股权比例低于已记录门槛 \\
行业不在范围内 & \textcolor{prismgreen}{绿色} & 所选行业未纳入审查机制 \\
无审查机制 & \textcolor{gray}{灰色} & 所选年份无有效正式机制 \\
\bottomrule
\end{tabularx}
\caption{交易筛查器结论等级}
\end{table}

结论附带通俗说明，内容涵盖：覆盖范围、门槛比较、申报义务、欧盟/非欧盟差异（如适用）、最长审查期限及主管机构。

\uielem{查看 [国家] 完整档案 →}按钮可立即跳转至国家标签页进行深入查阅。

\begin{warnbox}[法律免责声明]
交易筛查器提供基于 PRISM 数据库公开数据的指导性评估，不构成法律意见。在推进任何交易前，请务必咨询相关国家主管机构或具有资质的法律顾问。
\end{warnbox}

%% ============================================================
\section{AI 数据助手}
\label{sec:ai}

AI 助手支持以自然语言方式查询 PRISM 数据。

\begin{tipbox}[打开方式]
点击顶部导航栏中的\uielem{AI 助手}。
\end{tipbox}

\subsection{功能能力}

助手具有上下文感知能力：它了解当前选定的年份、国家、激活的分组筛选以及所在标签页。它可以：

\begin{itemize}
  \item 跨国比较审查机制（"2022 年德国与法国的审查机制有何差异？"）。
  \item 识别趋势（"哪些国家在 2019 年后收紧机制最多？"）。
  \item 解释数据字段（"净效益测试是什么意思？"）。
  \item 列出具有特定功能的国家（"2023 年哪些国家将 AI 行业纳入审查？"）。
  \item 总结具体法规（通过任意政策文件抽屉中的\uielem{用 AI 总结 →}按钮触发）。
\end{itemize}

\subsection{快捷提示词（芯片按钮）}

助手根据当前标签页和选定国家自动生成 3--4 个情境相关的提示词建议，点击任意芯片按钮即可立即发送。

\subsection{配置设置}

点击\uielem{⚙}图标可配置 AI 后端：

\begin{itemize}
  \item \textbf{服务提供商}——Anthropic（Claude）、OpenAI、DeepSeek、MiniMax、智谱 GLM、月之暗面 Kimi，或任意兼容 OpenAI 接口的自定义端点。
  \item \textbf{模型}——按提供商预填推荐模型，可手动覆盖。
  \item \textbf{基础 URL}——仅自定义提供商需要填写。
  \item \textbf{API 密钥}——存储于浏览器本地存储；仅直接发送至所选提供商的 API 端点，不经过任何其他传输。
\end{itemize}

\begin{tipbox}[键盘快捷键]
按 \keycap{Enter} 发送消息；在输入框中换行请使用 \keycap{Shift+Enter}。
\end{tipbox}

\subsection{局限性}

\begin{itemize}
  \item 助手的回答基于查询时注入的数据上下文，而非实时网络访问。
  \item 对于定量问题，助手接收严格程度前 8 名排名、所选国家完整档案以及查询中提及的行业覆盖列表。
  \item 回答字数上限约为 250 字（对比类列表可适当延长）。
  \item 回答质量取决于所选 AI 服务商和模型。
\end{itemize}

%% ============================================================
\section{全局搜索}
\label{sec:search}

顶部导航栏中的搜索框支持三种搜索模式：

\begin{table}[h]
\centering
\begin{tabularx}{\linewidth}{lX}
\toprule
\textbf{模式} & \textbf{搜索范围} \\
\midrule
国家 & 国家名称；点击结果可选定该国并打开国家标签页 \\
法规 & 法规标题及事件数据库中的国家名称 \\
政策文本 & 跨所有 141 份政策文件的全文检索；匹配段落高亮显示 \\
\bottomrule
\end{tabularx}
\caption{搜索模式}
\end{table}

\begin{itemize}
  \item 使用结果面板顶部的\uielem{国家 / 法规 / 政策}切换按钮切换搜索模式。
  \item 政策文本搜索结果在详情抽屉中显示，匹配段落以琥珀色高亮标注，并附下载或查看原始文件的链接。
  \item 按 \keycap{↑}/\keycap{↓} 浏览结果，按 \keycap{Enter} 选定，按 \keycap{Esc} 关闭。
\end{itemize}

%% ============================================================
\section{URL 分享}
\label{sec:url}

PRISM 自动将浏览器 URL 与当前状态保持同步：

\begin{verbatim}
https://jonebroom.github.io/prism-fdi/?tab=country&country=Germany&year=2021
\end{verbatim}

\begin{itemize}
  \item \texttt{tab}——当前激活的标签页名称。
  \item \texttt{country}——所选国家（使用数据中的标准英文名称）。
  \item \texttt{year}——所选年份。
\end{itemize}

随时从浏览器地址栏复制 URL，即可创建可分享链接，其他用户打开后将还原完全相同的视图。

%% ============================================================
\section{数据字段参考}
\label{sec:fields}

\subsection{定量字段}

\begin{table}[h]
\centering
\begin{tabularx}{\linewidth}{llX}
\toprule
\textbf{字段} & \textbf{类型} & \textbf{定义} \\
\midrule
\fieldname{num\_mechanisms} & 整数 & 当前有效的审查文件数量 \\
\fieldname{strictness} & 0--13 & 13 项二元程序性指标之和 \\
\fieldname{threshold} & \%（0--100） & 触发审查的最低股权收购比例 \\
\fieldname{time\_frame\_days} & 天 & 法定最长审查期 \\
\fieldname{total\_sectors} & 整数 & 明确涵盖的行业数量 \\
\bottomrule
\end{tabularx}
\caption{定量数据字段}
\end{table}

\subsection{13 项程序性权力指标}

每项均为二元变量（0/1），依据正式立法及监管指导文件编码：

\begin{table}[h]
\centering
\begin{tabularx}{\linewidth}{lX}
\toprule
\textbf{指标} & \textbf{含义} \\
\midrule
正式主动调查权 & 主管机构可自行启动审查，无需当事方申报 \\
增量股权审查 & 超过门槛后的后续股权增持也须接受审查 \\
申报费 & 申报时须缴纳强制性费用 \\
缓解措施 & 主管机构可施加条件（补救措施），而非全面禁止 \\
罚款 & 对违规或不合规行为可处以货币罚款 \\
国家安全测试 & 审查标准中含有明确的国家安全标准 \\
净效益测试 & 对东道国的经济效益是审查标准之一 \\
竞争测试 & 对市场竞争的影响是审查标准之一 \\
跨部门联合审查 & 多个政府机构参与审查程序 \\
分级管辖权 & 根据行业或交易规模，由不同机构行使管辖权 \\
当地代表 & 收购方须在东道国任命当地代表或董事会成员 \\
强化政府控制权 & 政府在批准后保留特殊股权或否决权 \\
联合选址 & 实体基础设施须在东道国境内联合选址 \\
\bottomrule
\end{tabularx}
\caption{13 项程序性权力指标}
\end{table}

\subsection{覆盖类型分类}

\begin{table}[h]
\centering
\begin{tabularx}{\linewidth}{lX}
\toprule
\textbf{类型} & \textbf{定义} \\
\midrule
行业性 & 审查仅限于明确列举的行业投资 \\
跨行业 & 所有行业均可能受到审查；覆盖范围最广 \\
混合型 & 结合强制性行业规则与酌情性更广泛审查 \\
资产型 & 由资产性质（如关键基础设施）触发，而非行业 \\
草案 & 立法已提出但尚未颁布 \\
已通过（未实施） & 法律已颁布；实施条例或生效日期尚未到位 \\
\bottomrule
\end{tabularx}
\caption{覆盖类型分类}
\end{table}

\subsection{通知与预批制度}

两个字段使用相同的术语体系：
\begin{itemize}
  \item \textbf{非强制}——无正式申报义务；仅可自愿提交。
  \item \textbf{强制（部分）}——符合特定条件（行业、门槛或投资者来源）的交易须申报。
  \item \textbf{强制（所有）}——所有超过最低限额的外国投资均须申报。
\end{itemize}

%% ============================================================
\section{常见问题解答}
\label{sec:faq}

\begin{enumerate}[label=\textbf{问\arabic*.}]

\item \textbf{为什么某国显示在地图上，却在国家标签页中显示"无审查机制"？}\\
地图显示全部 38 个国家。若某国在所选年份无机制，可能是：（a）所选年份早于该国建立机制的时间；或（b）该国从未建立正式审查制度。

\item \textbf{门槛字段显示"—"代表什么？}\\
表示以下情况之一：该国机制适用于所有外国投资而无具体股权门槛；门槛规定于未在此层级记录的附属法规中；或公开来源中无相关信息。

\item \textbf{为什么部分法规名称显示为"引用：34/2020"而非完整标题？}\\
部分来源数据集仅记录了法律引用编号（如法案编号和年份），而无完整标题。此类记录自动加注"引用："前缀，表明该名称为裸引用，而非描述性标题。

\item \textbf{严格程度评分如何计算？}\\
严格程度是 13 项二元程序性权力指标的简单求和。13 项权力全部激活的国家得 13 分；仅有国家安全测试和罚款的国家得 2 分。这是程序复杂性的衡量指标，而非审查导致禁止决定频率的指标。

\item \textbf{能否导出数据？}\\
平台目前不包含内置导出功能。基础数据集在平台数据说明中有所描述。如需用于研究，请联系数据集创建方。

\item \textbf{AI 助手表示无法回答我的问题，为什么？}\\
助手仅使用 PRISM 中的数据。超出 38 个国家、2007--2023 年范围的问题，或需要实时法律建议的问题，将收到免责说明。请确保 API 密钥和服务提供商配置正确。

\item \textbf{交易筛查器的输出是否具有法律效力？}\\
不具有法律效力。它提供基于编码数据的指导性自动评估。FDI 审查义务取决于交易具体事实、监管解释及频繁变化的政策指引。请务必向相关国家主管机构及具有资质的法律顾问核实。

\item \textbf{为什么地图上我的国家显示为灰色，即使分组筛选已包含该国？}\\
灰色阴影表示：（a）所选年份该国无审查机制；或（b）所选地图维度对该国-年份组合无数据（如门槛数据缺失）。将鼠标悬停于该国可查看说明数据可用情况的提示框。

\item \textbf{我分享的 URL 打开后与预期状态不同。}\\
URL 还原要求目标国家名称完全匹配。国家名称使用标准英文（如"United States"、"Republic of Korea"）。若 URL 创建时的年份超出数据范围（2007--2023），年份将被自动修正至最近有效值。

\end{enumerate}

%% ============================================================
\section{数据来源与方法论}
\label{sec:methodology}

\subsection{主要数据来源}

平台基于 \textbf{PRISM ISM 数据集}（2023.12 版）构建，该数据集是从政府官方来源、立法数据库及学术文献中整理而成的 FDI 审查机制结构化数据库。

\subsection{编码方法论}

\begin{itemize}
  \item 若某功能在现行立法中明确存在，二元指标编码为 \textbf{1}，否则编码为 \textbf{0}。模糊情形编码为 0（保守处理原则）。
  \item 门槛值反映存在分级门槛时最低适用门槛。
  \item 时限值反映初始阶段法定最长审查期（不含可能的延期）。
  \item 覆盖类型在多个文件并存时，反映主要现行机制。
\end{itemize}

\subsection{已知局限性}

\begin{itemize}
  \item 次国家层面及行业专项的更低门槛可能未被单一门槛字段完整捕捉。
  \item 应急和临时立法（包括新冠疫情措施）在已识别的情况下加以标注，但各司法管辖区的记录可能不一致。
  \item 少量历史记录（2007 年前）收录于有较长立法历史的国家（如芬兰 1939 年），但不属于主要分析时间序列。
  \item 政策文件全文来源于公开的 HTML 和 PDF 版本；PDF 提取过程中产生的格式痕迹可能出现在文本查看器中。
\end{itemize}

%% ============================================================
\section{快速参考卡片}
\label{sec:quickref}

\begin{tcolorbox}[enhanced, colback=lightgray, colframe=midgray, arc=6pt, title={\textbf{PRISM 功能速查}}, fonttitle=\bfseries]
\small
\begin{tabular}{@{}p{4.5cm}p{9cm}@{}}
\textbf{查看国家档案} & 点击侧边栏国家列表，或点击世界地图上的国家 → 查看详情 \\[3pt]
\textbf{按分组筛选} & 切换侧边栏中的 OECD / 欧盟/欧洲经济区 / 五眼联盟按钮 \\[3pt]
\textbf{更改年份} & 拖动年份滑块，或点击任意时间序列图中的年份 \\[3pt]
\textbf{时间动态演示} & 点击▶ 播放时间轴（切换标签页时停止） \\[3pt]
\textbf{核查交易} & ⚡ 交易核查 → 填入国家、行业、股权比例、来源国 \\[3pt]
\textbf{提问查询} & AI 助手 → 输入问题或点击芯片提示词 \\[3pt]
\textbf{查找法规} & 按 / → 切换至政策模式 → 输入关键词 \\[3pt]
\textbf{分享视图} & 从地址栏复制 URL \\[3pt]
\textbf{查看政策文本} & 点击任意法规气泡或事件卡片 → 政策文本显示于抽屉中 \\[3pt]
\textbf{跨年度比较} & 国家标签页 → 雷达图 → 选择"对比 [年份]"下拉菜单 \\[3pt]
\end{tabular}
\end{tcolorbox}

\vspace{1em}
\begin{tcolorbox}[enhanced, colback=prismclay!6, colframe=prismclay!40, arc=6pt, title={\textbf{数据字段速查}}, fonttitle=\bfseries]
\small
\begin{tabular}{@{}p{4cm}p{9.5cm}@{}}
\textbf{严格程度（0--13）} & 13 项二元程序性权力之和。数值越高，机制越复杂。 \\[3pt]
\textbf{门槛（\%）} & 触发审查的最低股权比例（空白 = 适用于所有 / 未知）。 \\[3pt]
\textbf{时限（天）} & 初始阶段法定最长审查期。 \\[3pt]
\textbf{覆盖类型} & 行业性 → 跨行业 → 混合型 → 资产型（覆盖范围递增）。 \\[3pt]
\textbf{预批} & 须在交割\textit{前}取得许可。 \\[3pt]
\textbf{通知} & 须向主管机构通知（决定作出前可完成交割）。 \\[3pt]
\end{tabular}
\end{tcolorbox}

\vspace{2em}
\hrule
\vspace{0.5em}
{\small\color{gray}
\textbf{PRISM ISM 数据库} · \url{https://jonebroom.github.io/prism-fdi/}\\
数据来源：PRISM ISM 数据集 v2023.12 · 38 个国家 · 2007--2023\\
本手册依据平台源代码与数据结构自动生成。
}

\end{document}
